\documentclass[12pt]{article}
\usepackage[utf8]{inputenc}
\usepackage[T1]{fontenc}
\usepackage{hyperref}
\title{Phupa Language Resources for K-12 Curriculum Development}
\author{Compiled by Claude Sonnet 5.5 (claude-sonnet-5-5) for Andr\'as Kornai}
\date{}
\begin{document}
\maketitle

\section{Introduction}
Phupa (ISO 639-3 ypp; Glottocode phup1239) is a Loloish (Sino-Tibetan) language of Yunnan. It is spoken in southwest Mengzi County and in southeast Lengquan and Shuitian townships on the southeast Gejiu panhandle. Joshua Project gives 3,300 speakers (Ethnologue 2016). Wikipedia describes the Phupha dialect cluster, of which Phupa is one variety, at 14,000 speakers (2007). The two figures cover different scopes. Wikipedia says use is decreasing, and Joshua Project notes growing use of Mandarin among the young. The same Wikipedia page says that the Alugu variety, which is not Phupa, is taught in primary schools. Nothing found shows Phupa itself being taught.

\section{Primary Educational Resources}
No primers, textbooks or school programmes for Phupa were found.

\section{Linguistic and Reference Materials}
\begin{itemize}
\item Wikipedia, Phupha language (the cluster that includes Phupa): \url{https://en.wikipedia.org/wiki/Phupha_language}
\item Wikipedia mirror, ISO 639:ypz: \url{https://wikipedia.com/wiki/ISO_639:ypz}
\end{itemize}

\section{Religious and Community Resources}
Joshua Project pages. They are written from a missionary perspective:
\begin{itemize}
\item Phupa: \url{https://joshuaproject.net/people_groups/20611}
\item Phupha, a related group: \url{https://joshuaproject.net/people_groups/20628}
\end{itemize}

\section{Digital and Audio Resources}
None were found.

\section{Curriculum Development Context}
Published material is minimal. The only schooling mentioned in the results concerns the related Alugu variety, so it is a possible lead to follow up and not a Phupa resource. Further work would need documentation and fieldwork together with the community.
\end{document}